\scanpage{14}
\mspage{28}

1) $B$ est formellement simple sur $A$

2) $B$ satisfait la condition P \quad 3) $B$ est plat sur $A$ \quad 3 bis) Pour tout
entier $n\geq0$, $B_n=B/I^{n+1}B$ est plat sur $A_n=A/I^{n+1}$. \deleted{4)}

Il reste seulement à montrer l'équivalence de 3) et 3 bis), qui est
contenue en effet (moyennant les hypothèses noethériennes faites) dans
la platitude supérieure.

\deleted{Le résultat précédent} Il conviendrait de noter aussi \inserted{dès après 3.7.} que ``algèbre de
Cohen'' est une notion stable par extension de la base via produits tensoriels complétés: c'est une trivialité.

Le résultat \inserted{4.}8. montre l'intérêt qu'il y a, $A,I$ et $B_0$ étant donnés,
de disposer d'un \uline{théorème d'existence} pour une algèbre de Cohen $B$ sur $A$
se réduisant suivant $B_0$. On notera que la construction habituelle de
proche en proche nous ramène \deleted{au cas où $A$ est discr}(remplaçant $A$ par des $A_i$
et $I$ par des puissances) au cas où $I$ est de carré nul, (donc où $A$ est
discret). On obtient alors un problème de nature essentiellement infinitésimale,
dont nous ignorons s'il existe toujours une solution, faute de
savoir caractériser la structure des algèbres de Cohen $B_0$ sur l'anneau
discret $A_0$. Cependant, nous obtenons \inserted{(4.10.)} un théorème d'existence dans
un cas particulier intéressant, en utilisant les résultats du N°2. Notons
d'abord

\deleted{\uline{Théorème} \deleted{3.9.} Soit $A$ \deleted{$B$ un homomorphisme} un anneau local noethérien,
\deleted{et de corps} résiduel $k$, et soit $B_0$ une $k$-algèbre,
\deleted{qui soit un satisfasse}
\deleted{qui soit}
\deleted{aux conditions suivantes: $B_0$ est} un anneau local noethérien, \deleted{et complet,}
\inserted{de} \deleted{son corps} résiduel $K$
\deleted{est de degré de multiplicité radicielle finie sur $k$}
\deleted{(ce qui est le cas} une extension séparable de $k$, $B_0$ est (pour sa topolo-}

\uline{Corollaire} \inserted{4.}9. Soit $A\longrightarrow B$ un homomorphisme local d'anneaux locaux noethériens,
\deleted{tel que l'extension résiduelle $k(B)/k(A)$ soit à multiplicité}

\mspage{13}
